%Author K.FUJIWARA
\documentclass{bxjsarticle}
%[[[ TeX setting
%[[[ package
\usepackage{zxjatype}
\usepackage{amsmath,amsthm}
\usepackage{unicode-math}
\usepackage{fontspec}
\setCJKmainfont{MogaExMincho}
\setmathfont{NewComputerModernMath}
\usepackage[xetex]{graphicx}
\usepackage[xetex]{animate}
%]]]

%[[[ env setting
\newtheorem{Theorem}{Theorem}%[section]
\newtheorem{Definition}[Theorem]{Definition}
\newtheorem{Corollary}[Theorem]{Corollary}
\newtheorem{Lemma}[Theorem]{Lemma}
\newtheorem{Proposition}[Theorem]{Proposition}
\newtheorem{Remark}{Remark}%[section]
%\addtolength{\textwidth}{5.cm}
%\addtolength{\textheight}{4cm}
%]]]
%]]]

%[[[ doc
\begin{document}

\title{全微分可能性の補足}
\author{藤原和将}
\date{2021年11月10日}
\maketitle

%[[[ Intro
$D ⊂ ℝ^2$ ($D ∋ \{0\}$)を開集合とする.
$f: D → ℝ$が原点に於いて全微分可能であるとは,
$f$が２つの定数$A, B$を伴って,
	\[
	f(x,y)
	=Ax + By + o(|(x,y)|)
	\]
を満たす事である.
このノートでは,
原点に於ける全微分可能性の判定方法に就いて見る.

$f$が原点で全微分可能である事を判定する方法として,
次がよく用いられる.

%[[[ Proposition : Hantei
\begin{Proposition}
\label{P:C1}
$f ∈ C^1(D)$ならば$f$は原点で全微分可能である.
特に,
	\[
	A = ∂_x f(0,0),\quad
	B = ∂_y f(0,0)
	\]
である.
\end{Proposition}
%]]]

一方で, $f$の極座標表示が特定の形をしているならば,
原点での全微分可能性の判定は容易である.

%[[[ Proposition : kyokuzahyou
\begin{Proposition}
\label{P:Rad}
$φ: [-π,π] → ℝ$が$φ(-π)=φ(π)$を満たすとする.
$f : ℝ^2 → ℝ$を
	\[
	f(r \cos θ, r \sin θ) = r φ (θ)
	\]
で与える.
この時, $f$が原点に於いて全微分可能である事と
	\begin{align}
	f(x,y) = A x + B y
	\label{eq:1}
	\end{align}
なる実数$A,B$が存在する事は同値である.
\end{Proposition}

%[[[ Remark : r φ
\begin{Remark}
上の命題を言い換えれば,
$f$が原点に於いて全微分可能である事と, $φ(θ) = A \cos θ + B \sin θ$である事が同値である.
また, $f$が原点に於いて全微分可能である事と,
$f$のグラフが原点を含む平面となることが同値である.
特に, 原点で全微分可能であるならば,
$f$は$ℝ^2$上至る所全微分可能である.
\end{Remark}
%]]]

\begin{proof}
\eqref{eq:1}ならば,
全微分可能性の定義より$f$は全微分可能である.
$f$が全微分可能であるならば,
２つの実数$A,B$が存在して,
	\[
	f(r \cos θ, r \sin θ) - f(0,0)
	= ( A \cos θ + B \sin θ ) r + o(r)
	\]
である.
$f(0,0)=0$なので,
	\begin{align*}
	φ(θ)
	&= \lim_{r ↘ 0} \frac{f(r \cos θ, r \sin θ)-f(0,0)}{r}\\
	&= \lim_{r ↘ 0} \frac{( A \cos θ + B \sin θ ) r + o(r)}{r}\\
	&= A \cos θ + B \sin θ
	\end{align*}
であるから,
$o(r)=0$である.
従って, $f(x,y)=Ax+By$を得る.
\end{proof}
%]]]

%[[[ Remark : 例
\begin{Remark}
上記命題より, 例えば,
	\[
	b(x,y) =
	\begin{cases}
	\frac{x^3+y^3}{x^2+y^2}
	&\mathrm{if} \quad (x,y) ≠ 0,\\
	0
	&\mathrm{if} \quad (x,y) = 0
	\end{cases}
	\]
は原点に於いて連続且つ偏微分可能であるが, 全微分不能である.
実際
	\[
	b(r \cos θ, r \sin θ)
	= r ( \cos^3 θ + \sin^3 θ)
	\]
である.
$b(x,y)$は一次式ではないので,
$b$は原点に於いて全微分不能である.
偏導関数が原点で不連続になっている事は,
Proposition \ref{P:C1}の対偶命題から従う.
また, Remark 4の計算から直接見て取れる.
\end{Remark}
%]]]

%[[[ Remark
\begin{Remark}
$F$を$f$の極座標表示とすれば,
	\begin{align*}
	\frac{∂}{∂x} f(0,0)
	&= \frac{∂}{∂ r} F(0,0) = φ(0),\\
	\frac{∂}{∂y} f(0,0)
	&= \frac{∂}{∂ r} F(0,π/2) = φ(π/2)
	\end{align*}
である.
\end{Remark}
%]]]

%[[[ Remark
\begin{Remark}
上記命題で$φ$が微分可能な場合に於いて,
$f$の偏導関数の挙動を見てみよう.
$F$を$f$の極座標表示とすれば,
$(x,y) ≠ (0,0)$に於いて
	\begin{align*}
	\frac{∂}{∂x} f(x,y)
	&= \cos θ \frac{∂}{∂ r} F(r,θ) - \frac{\sin θ}{r} \frac{∂}{∂ θ} F(r,θ)\\
	&= φ(θ) \cos θ - φ'(θ) \sin θ,\\
	\frac{∂}{∂y} f(x,y)
	&= \sin θ \frac{∂}{∂ r} F(r,θ) + \frac{\cos θ}{r} \frac{∂}{∂ θ} F(r,θ)\\
	&= φ(θ) \sin θ + φ'(θ) \cos θ
	\end{align*}
である.
もし, 上記２つの偏導関数が原点に於いて連続であるならば,
２つの偏導関数は$θ$に依存しない.
即ち,
	\[
	φ(θ) \cos θ - φ'(θ) \sin θ = A,\quad
	φ(θ) \sin θ + φ'(θ) \cos θ = B
	\]
なる定数$A,B$が存在する.
	\[
	φ(θ)
	= \cos θ \big\{ φ(θ) \cos θ - φ'(θ) \sin θ \}
	+ \sin θ \big\{ φ(θ) \sin θ + φ'(θ) \cos θ \}
	= A \cos θ + B \sin θ
	\]
より, Proposition \ref{P:C1}とProposition \ref{P:Rad}の関係が分かる.
\end{Remark}
%]]]

\noindent
{\bf 参考}\\
鈴木 武, 柴田 良弘, 田中 和永 , 山田 義雄 著,
理工系のための微分積分I,
2007/4/1 内田老鶴圃出版
(ISBN: 4753601811)


\end{document}
%]]]
